\documentclass[12pt]{article}
\usepackage[a4paper, margin=1in]{geometry}
\usepackage{graphicx}
\usepackage{titlesec}
\usepackage{fancyhdr}
\usepackage{hyperref}
\usepackage{times}
\usepackage{array}
\usepackage{longtable}
\usepackage{booktabs}
\usepackage{enumitem}
\usepackage{microtype}
\usepackage{multirow}
\usepackage{float}
\usepackage{caption}

% Page Style Setup
\pagestyle{fancy}
\fancyhf{}
\cfoot{\thepage}
\renewcommand{\headrulewidth}{0pt}
\renewcommand{\footrulewidth}{0pt}

% Section Title Formatting
\titleformat*{\section}{\Large\bfseries\centering}
\titleformat*{\subsection}{\large\bfseries}
\titlespacing*{\section}{0pt}{3.5ex plus 1ex minus .2ex}{2.3ex plus .2ex}
\titlespacing*{\subsection}{0pt}{3.25ex plus 1ex minus .2ex}{1.5ex plus .2ex}

\begin{document}

% Title Page
\begin{titlepage}
    \centering
    \vspace*{1cm}
    
    {\Large\bfseries TrackEneer: A Smart Academic Companion\\[0.2cm]}
    
    \vspace{1cm}
    
    {\large Submitted in partial fulfilment of the requirements of\\[0.5cm]
    MAJOR PROJECT\\[0.3cm]
    SEMESTER VII \& VIII\\[0.5cm]
    for the degree of\\[0.5cm]
    \bfseries BACHELOR OF ENGINEERING\\[0.2cm]
    in\\[0.2cm]
    \bfseries COMPUTER SCIENCE AND ENGINEERING
    (Cybersecurity,IoT including Blockchain Technology)}\\[0.5cm]

    {\large By} \\[0.5cm]

    \begin{tabular}{|c|l|l|}
        \hline
        \textbf{Sr. No.} & \textbf{Name}     & \textbf{Student ID} \\ \hline
        1                & Himanshu Pandey   & 202204043           \\ \hline
        2                & Kulsum Khan       & 2023042004          \\ \hline
        3                & Bhanavi Pandey    & 202204042           \\ \hline
        4                & Uday Naik         & 202204041           \\ \hline
    \end{tabular}

    \vspace{1cm}

    {\large under the\\[0.3cm] guidance of} \\[0.5cm]

    {\textbf {\large Prof. Suhas Lawand}}\\[0.1cm]
    
    \vspace{1cm}
    
    {\large \textbf{Xavier Institute of Engineering}}
    \vspace{1cm}

    {\large Department of Computer Engineering,\\[0.2cm]
    Xavier Institute of Engineering, Mahim,\\[0.2cm]
    Mumbai - 400 016}\\[0.5cm]

    {\large University of Mumbai}\\[0.3cm]

    {\large (AY 2024-25)}
\end{titlepage}

% Certificate Page
\newpage
\section*{CERTIFICATE}
\vspace{1.5cm}
\noindent
This is to certify that the project entitled \textbf{``TrackEneer: A Smart Academic Companion''} is a bona-fide work of Himanshu Pandey, Kulsum Khan, Bhanavi Pandey, and Uday Naik submitted to the University of Mumbai in partial fulfillment of the requirement for the award of the degree of \textbf{``Bachelor of Engineering''} in \textbf{``Computer Science and Engineering(Cybersecurity,IoT including Blockchain Technology)''}.

\vspace{4cm}

    \parbox{6cm}{
        \centering
        \rule{6cm}{0.4pt} \\
        \vspace{0.2cm}
        \textbf{Prof. Suhas Lawand} \\
        Supervisor
    }

\thispagestyle{empty}

% Approval Page
\newpage
\section*{PROJECT APPROVAL}
\vspace{1.5cm}
\noindent
This Project entitled \textbf{``TrackEneer: A Smart Academic Companion''} by Himanshu Pandey, Kulsum Khan, Bhanavi Pandey, and Uday Naik is approved for the degree of \textbf{``Bachelor of Engineering''} in \textbf{``Computer Engineering''}.

\vspace{5cm}
\noindent\begin{parbox}[t]{\textwidth}{
    \raggedleft
    \parbox{6cm}{
        \centering
        \rule{6cm}{0.4pt} \\
        \vspace{0.2cm}
        \textbf{Prof. Suhas Lawand} \\
        Supervisor
    }
}\end{parbox}

\vspace{3cm}
\noindent\begin{parbox}[t]{\textwidth}{
    \raggedleft
    \parbox{6cm}{
        \centering
        \rule{6cm}{0.4pt} \\
        \vspace{0.2cm}
        \textbf{(External Examiner)}
    }
}\end{parbox}

\vfill
\noindent
\textbf{Date:} \par
\vspace{0.5cm}
\noindent
\textbf{Place: Mumbai}
\thispagestyle{empty} 

% Declaration Page
\newpage
\section*{Declaration}
\vspace{1.5cm}
\noindent
We declare that this written submission represents our ideas in our own words and where others' ideas or words have been included, we have adequately cited and referenced the original sources. We also declare that we have adhered to all principles of academic honesty and integrity and have not misrepresented or fabricated or falsified any idea/data/fact/source in our submission. We understand that any violation of the above will be cause for disciplinary action by the Institute and can also evoke penal action from the sources which have thus not been properly cited or from whom proper permission has not been taken when needed.

\vspace{3cm}
\begin{flushright}
    \parbox{7cm}{
    \centering
    \vspace{1cm}
    \rule{7cm}{0.4pt} \\
    \vspace{0.1cm}
    Himanshu Pandey (202204043)\\[1.5cm]

    \vspace{1cm}
    \rule{7cm}{0.4pt} \\
    \vspace{0.1cm}
    Kulsum Khan (2023042004)\\[1.5cm]

    \vspace{1cm}
    \rule{7cm}{0.4pt} \\
    \vspace{0.1cm}
    Bhanavi Pandey (202204042)\\[1.5cm]
    
    \vspace{1cm}
    \rule{7cm}{0.4pt} \\
    \vspace{0.1cm}
    Uday Naik (202204041)\\[1.5cm]
    }
\end{flushright}

\vfill
\noindent
\textbf{Date:}
\vspace{4cm}
\thispagestyle{empty}

% Table of Contents
\newpage
\section*{Contents}
\begin{center}
\setlength{\extrarowheight}{0.5cm} 
\begin{tabular}{|>{\raggedright\arraybackslash}p{0.8\textwidth}|>{\raggedright\arraybackslash}p{0.1\textwidth}|} 
\hline 
\textbf{Abstract} & i \\ \hline
\textbf{List of Figures} & ii \\ \hline
\textbf{List of Tables} & iii \\ \hline
\textbf{1. Introduction} & \textbf{1}  \\ \hline
\textbf{2. Literature Survey} & \textbf{3}\\ \hline
\textbf{3. Problem Statement and Objectives} & \textbf{4}\\ \hline
\textbf{4. Proposed System} & \textbf{5}\\ \hline
\textbf{5. Methodology (Proposed System)} & \textbf{7} \\ \hline
\textbf{6. Experimental Setup} & \textbf{9} \\ \hline
\textbf{7. Implementation Plan for Next Semester} & \textbf{12} \\ \hline
\textbf{8. Timeline Chart} & \textbf{13} \\ \hline
\textbf{9. Conclusion} & \textbf{14} \\ \hline
\textbf{References} & \textbf{15} \\ \hline
\end{tabular}
\end{center}
\thispagestyle{empty}

% Abstract
\newpage
\pagenumbering{roman}
\setcounter{page}{1}
\section*{Abstract}
\addcontentsline{toc}{section}{Abstract}
TrackEneer is an integrated AI-powered academic companion designed to holistically address the multifaceted challenges faced by engineering students throughout their academic journey. The system comprises four interconnected modules: Schedule, Study, Placement, and Insights, each leveraging modern artificial intelligence and web technologies to provide intelligent automation and personalization. The Schedule module implements an intelligent task management system with deadline tracking and live notifications. The Study module serves as a centralized repository with subject-wise organization, supporting multiple file formats and offering in-browser preview capabilities. The Placement module utilizes Google's Gemini AI to generate comprehensive company research reports, including vision analysis, placement process details, and personalized recommendations. The Insights module employs Cohere AI to deliver year and branch-specific academic guidance, providing students with customized roadmaps encompassing focus areas, study strategies, technical skills development, and career pathways. Built on a modern technology stack featuring FastAPI backends, Next.js frontend, and OAuth authentication, TrackEneer demonstrates the potential of hybrid intelligence systems in education, where artificial intelligence augments human cognitive capabilities to enhance academic productivity, reduce information fragmentation, and provide context-aware personalized learning experiences.

\newpage
\section*{List of Figures}
\addcontentsline{toc}{section}{List of Figures}
\listoffigures

\newpage
\section*{List of Tables}
\addcontentsline{toc}{section}{List of Tables}
\listoftables

% Main Content
\newpage
\pagenumbering{arabic}
\setcounter{page}{1}

\section{Introduction}
\subsection{Motivation and Background}
The contemporary engineering curriculum imposes substantial cognitive loads on students, who must navigate a complex ecosystem of academic materials, from dense syllabi and extensive lecture notes to the rigorous demands of placement preparation. While the proliferation of digital tools has offered some relief, it has paradoxically led to a fragmented technological landscape. Students often find themselves juggling disparate platforms for note-taking, task scheduling, and career preparation, a practice that engenders inefficiency and exacerbates cognitive overload. A fundamental limitation of these conventional tools is their inability to semantically interpret academic documents, necessitating laborious manual data entry for tasks like creating study plans or organizing notes. This absence of intelligent automation constitutes a significant impediment to effective time and information management. This vision aligns with the emerging concept of hybrid intelligence systems in education, where the symbiotic interplay between human cognition and artificial intelligence serves not merely to automate tasks, but to augment and enhance the user's cognitive and learning capabilities \cite{b10}.

\subsection{Objectives}
The main objective of TrackEneer is to develop a smart academic companion that addresses the aforementioned problems. The specific objectives are:
\begin{enumerate}
    \item \textbf{Implement Intelligent Task Management:} To develop a schedule management system with deadline tracking, priority management, and automated notifications.
    \item \textbf{Create Centralized Note Repository:} To build a subject-organized study materials platform supporting multiple file formats with preview and download capabilities.
    \item \textbf{Enable AI-Powered Company Research:} To implement an intelligent placement preparation module that generates comprehensive company insights using generative AI.
    \item \textbf{Deliver Personalized Academic Guidance:} To develop an insights system that provides year and branch-specific academic roadmaps and recommendations.
    \item \textbf{Ensure Seamless Integration:} To create a unified dashboard connecting all modules with consistent user experience and OAuth authentication.
\end{enumerate}

\subsection{Outline of the Report}
This report is organized as follows: Section 2 reviews existing literature and identifies research gaps. Section 3 presents the problem statement and objectives. Section 4 describes the proposed TrackEneer system, including its framework and design. Section 5 details the methodology and implementation approach. Section 6 explains the experimental setup, data sources, and evaluation parameters. Section 7 outlines the implementation plan for the next semester. Section 8 provides a project timeline. Finally, Section 9 concludes the report.

\section{Literature Survey}
\subsection{Existing Systems}
The theoretical foundations of TrackEneer are rooted in several intersecting domains of research: educational technology, document intelligence, and automated scheduling. Existing literature provides a strong precedent for the individual components of our system, yet reveals a gap in their holistic integration.

In the field of document analysis, researchers have made significant strides. The work of Sodhi and Choudhary (2023), for example, demonstrated a system for simplifying syllabus creation through content extraction, affirming the viability of automating this preliminary academic task \cite{b1}. Similarly, AI-powered study planners, such as the one proposed by Lingshetti et al. (2023), have shown the potential of intelligent scheduling to create personalized learning pathways, a concept central to TrackEneer \cite{b4}.

This fragmentation forces students to context-switch between multiple applications, leading to inefficiencies. TrackEneer is designed to bridge this gap by creating a unified ecosystem where these functions are deeply interconnected, addressing the limitations of prior work.

\subsection{Theoretical Foundations}
TrackEneer's design draws on several core theoretical ideas spanning information retrieval, representation learning, graph theory, scheduling optimization, and human-in-the-loop AI. Below we briefly introduce each area and the rationale for its use within the system.

\paragraph{Semantic Representations and Embeddings}
Modern information retrieval benefits from dense vector representations (embeddings) that capture semantic similarity beyond lexical overlap. We use Sentence-Transformers ("all-mpnet-base-v2") to map documents, syllabus items, and short answers into a shared embedding space. In this space, semantic search reduces to nearest-neighbor queries in high-dimensional vector spaces, where the cosine similarity approximates semantic relatedness. Key theoretical properties exploited are the continuity of the embedding function and the preservation of semantic neighborhoods under the learned mapping.

\paragraph{Approximate Nearest Neighbors and Vector Stores}
Exact nearest-neighbor search in high-dimensional spaces is computationally expensive; TrackEneer uses an ANN-backed vector store (Chroma/FAISS-like indexing) to perform sub-linear time retrieval. The trade-off between retrieval latency and recall is managed via index settings (number of probes, PQ resolution). The system caches frequently-used vectors and query results to reduce repeated I/O and indexing overhead, improving end-to-end latency for interactive tasks.

\paragraph{Graph Representations for Scheduling}
Course pre-requisites, topic dependencies, and temporal constraints are naturally modeled as directed graphs. We use a Neo4j-backed graph to represent dependency structures and to compute constrained topological orderings for study plans. Scheduling reduces to solving constrained ordering and resource allocation problems (e.g., time budgets, deadlines), for which heuristics and shortest-path formulations provide efficient, practical solutions even when optimal combinatorial solutions are intractable.

\paragraph{Probabilistic Ranking and Re-ranking}
Initial candidate retrieval via embeddings is often followed by learned or heuristic re-ranking that incorporates signals such as recency, relevance feedback, user performance (weak labels), and domain-specific importance. TrackEneer supports lightweight re-ranking layers (Cohere/Gemini scoring or logistic mixtures) to better match user intent to final presented items.

\paragraph{Hybrid Intelligence and Human-in-the-Loop}
To maintain high precision for critical tasks (e.g., placement feedback), the architecture supports hybrid workflows where automated recommendations are presented alongside confidence scores and interfaces for human validation. This reduces catastrophic failures from model hallucination and enables continual improvement via labeled user corrections.

\paragraph{Real-time Interaction and Event Streams}
Interactive scheduling and live notifications rely on low-latency event streams. WebSocket-backed channels (notification service) provide near real-time delivery of schedule updates and reminders. The theoretical model abstracts this as an event-driven system with eventual consistency guarantees for cached views.

These foundations inform the system architecture choices documented throughout the report and motivate the hybrid use of dense retrieval, graph reasoning, and real-time pipelines.

\section{Problem Statement and Objectives}
\subsection{Problem Statement}
Engineering students currently lack an intelligent, centralized system that can seamlessly integrate their academic and placement preparation activities. The primary challenges include:
\begin{itemize}
    \item \textbf{Manual Academic Management:} Students manually process syllabi, organize notes, and create study schedules, which is time-consuming and prone to error.
    \item \textbf{Fragmented Resources:} Study materials, notes, and placement resources are often scattered across multiple platforms, leading to disorganization.
    \item \textbf{Lack of Personalization:} Generic study planners and placement portals do not adapt to individual learning styles, strengths, or career goals.
    \item \textbf{Inefficient Placement Preparation:} Curating relevant placement materials like past papers and mock tests is a manual and often overwhelming task.
    \item \textbf{Poor Time Management:} Students struggle with deadline tracking and prioritization across multiple subjects and activities.
\end{itemize}

\subsection{Objectives}
The objectives have been previously outlined in Section 1.2.

\section{Proposed System}
\subsection{System Overview}
TrackEneer is proposed as an integrated, AI-driven solution to address the challenges of fragmented academic management. It serves as a unified platform comprising four core modules:

\begin{enumerate}
    \item \textbf{Schedule Module:} Intelligent task and deadline management with priority tracking and live notifications.
    \item \textbf{Study Module:} Centralized note repository with subject-wise organization, multi-format support, and preview capabilities.
    \item \textbf{Placement Module:} AI-powered company research tool generating comprehensive insights using Google Gemini AI.
    \item \textbf{Insights Module:} Personalized academic guidance system providing year and branch-specific roadmaps using Cohere AI.
\end{enumerate}

These modules are unified through a central dashboard with OAuth authentication, ensuring seamless navigation and consistent user experience across all functionalities.

\subsection{Module Flowcharts}
The following flowcharts summarize the high-level interactions and data flow for each core module. The PNG files referenced here are included by filename so they can be uploaded directly to Overleaf.
\begin{figure}[htbp]
    \centering
    \includegraphics[width=\textwidth]{architecture_system_diagram.png}
    \caption{System Architecture}
    \label{fig:architecture_system}
\end{figure}


\subsection{Analysis/Framework/Algorithm}
The architecture of TrackEneer is conceived as a multi-layered system to ensure modularity, scalability, and the efficient orchestration of data and intelligence. It comprises three primary strata: a presentation layer (frontend), a service layer (backend), and an intelligence layer (AI models).

\subsubsection{Frontend}
The frontend, developed using \textbf{Next.js 15} with \textbf{React 19} and styled with \textbf{Tailwind CSS 4}, serves as the human-computer interface. Its design philosophy prioritizes a responsive and intuitive user experience with a consistent dark theme across all modules. The application implements file-based routing for seamless navigation between the dashboard, schedule, study, placement, and insights modules.

\subsubsection{Backend}
The backend implements a unified FastAPI application with modular routers, providing logical separation of concerns while simplifying deployment and maintenance. Important endpoints for the main modules are exposed via routers mounted on the unified backend (port 5000). The notification service remains a separate Express/TypeScript server for Web Push integration and runs on port 3001.

Each router (Scheduler, Study, Placement, Insights) uses JSON-based cache files during development and can be migrated to persistent databases (e.g., MongoDB, Neo4j) for production. \textbf{NextAuth.js} handles OAuth authentication with Google, providing secure user sessions across the frontend and backend.

\subsubsection{AI Layer}
The AI layer represents the core intellectual capital of TrackEneer, utilizing state-of-the-art language models:
\begin{itemize}
    \item \textbf{Google Gemini 2.5 Flash:} Employed in the placement router for generating comprehensive company insights, analyzing vision/mission statements, placement processes, and providing personalized recommendations. The router constructs structured prompts, parses responses, and leverages a 24-hour JSON cache to minimize repeated API calls and latency.

    \item \textbf{Cohere Command-R (08-2024):} Utilized in the insights router for creating year and branch-specific academic guidance. The router generates detailed preparation roadmaps covering focus areas, study strategies, technical skills, project recommendations, placement timelines, career guidance, extracurricular activities, and common pitfalls to avoid.

    \item \textbf{Caching Strategy:} Both AI-powered routers implement a 24-hour cache system (JSON files in development) that stores generated content, reducing API costs and providing fast responses for commonly requested items.
\end{itemize}

\subsection{Design Details}
The system is designed with a user-centric philosophy emphasizing:

\textbf{Consistent UI/UX:} All modules share a dark theme with cyan-blue-purple gradient design system. The minimalist interface reduces cognitive load while maintaining visual appeal and professional aesthetics.

	extbf{Modular Architecture:} Each backend router operates logically within the unified FastAPI backend, providing dedicated API endpoints that allow for independent development, maintenance, and replacement without affecting other components.

\textbf{Data Persistence:} JSON-based file storage provides simplicity during development while maintaining flexibility for future migration to database systems like MongoDB or PostgreSQL.

\textbf{Security:} OAuth 2.0 authentication eliminates password storage risks, while HTTPS enforcement and CORS policies protect against common web vulnerabilities.

\textbf{Responsive Design:} Tailwind CSS breakpoints ensure optimal viewing on mobile (< 640px), tablet (640px - 1024px), and desktop (> 1024px) devices.

\section{Methodology (Proposed System)}
\subsection{Schedule Module Implementation}
The schedule module implements a comprehensive task management system with the following workflow:

\begin{enumerate}
    \item \textbf{Task Creation:} Users create tasks with attributes including title, description, due date, priority level (High/Medium/Low), and subject association.
    
    \item \textbf{Storage and Retrieval:} Task data is persisted in \texttt{schedule\_data.json} with automatic indexing by date and priority for efficient retrieval.
    
    \item \textbf{Dashboard Integration:} The Next Tasks component displays upcoming deadlines with priority badges, color-coded for visual distinction.
    
    \item \textbf{Notification System:} Live notifications alert users to tasks due within 24 hours, implementing a time-based priority queue.
\end{enumerate}

\subsection{Study Module Implementation}
The study module provides centralized file management through the following process:

\begin{enumerate}
    \item \textbf{File Upload:} Users select a subject and upload files (PDF, DOCX, PNG, JPG, GIF, TXT). The backend validates file types and sizes before storage.
    
    \item \textbf{Metadata Generation:} The system automatically extracts and stores metadata including filename, subject, upload date, file size, and MIME type in \texttt{study\_notes.json}.
    
    \item \textbf{File Organization:} Files are stored in an \texttt{uploads/} directory with unique identifiers, preventing naming conflicts and enabling efficient retrieval.
    
    \item \textbf{Preview System:} For supported formats (PDF, images), the \texttt{/api/notes/\{id\}/preview} endpoint serves files with \texttt{Content-Disposition: inline}, enabling browser-based preview in a modal window.
    
    \item \textbf{Download Functionality:} The \texttt{/api/notes/\{id\}/download} endpoint serves files with \texttt{Content-Disposition: attachment}, preserving original filenames.
\end{enumerate}

\subsection{Placement Module Implementation}
The placement module leverages Google Gemini AI through the following pipeline:

\begin{enumerate}
    \item \textbf{Company Query:} Users input a company name and optionally a website URL.
    
    \item \textbf{Cache Check:} The system queries \texttt{placement\_cache.json} for existing data less than 24 hours old, returning cached results if available.
    
    \item \textbf{AI Generation:} If no cache exists, the system constructs a detailed prompt instructing Gemini to generate:
    \begin{itemize}
        \item Vision and Mission analysis
        \item Placement Process details (eligibility, rounds, timeline)
        \item Company Review (work culture, growth opportunities, compensation)
        \item Overall Score (1-10) with justification
        \item Strengths and Considerations lists
        \item Personalized Recommendations
    \end{itemize}
    
    \item \textbf{Response Processing:} The AI response is parsed, structured into JSON format, cached for 24 hours, and returned to the frontend.
    
    \item \textbf{Alumni Display:} Xavier Institute of Engineering alumni placement data is displayed separately, providing real-world success stories and salary benchmarks.
\end{enumerate}

\subsection{Insights Module Implementation}
The insights module delivers personalized guidance through:

\begin{enumerate}
    \item \textbf{Input Collection:} Users select their academic year (FE/SE/TE/BE) and optionally their branch (Computer Engineering, IT, Electronics, etc.).
    
    \item \textbf{Cache Verification:} Similar to the placement module, the system checks \texttt{insights\_cache.json} for existing guidance matching the year-branch combination.
    
    \item \textbf{Prompt Engineering:} A comprehensive prompt instructs Cohere Command-R to generate eight distinct sections:
    \begin{itemize}
        \item Focus Areas specific to the year and branch
        \item Study Strategy and Time Management recommendations
        \item Technical Skills roadmap with learning priorities
        \item Project Ideas categorized by difficulty (beginner/intermediate/advanced)
        \item Placement Preparation Timeline with month-by-month milestones
        \item Extracurricular Activities suggestions
        \item Career Guidance with potential career paths
        \item Common Mistakes to Avoid based on year-specific challenges
    \end{itemize}
    
    \item \textbf{Content Formatting:} The generated content is formatted with markdown parsing to properly render bold text, headers, and bullet points in the frontend.
    
    \item \textbf{Caching:} The complete guidance is cached for 24 hours, enabling instant retrieval for subsequent requests.
\end{enumerate}

\section{Experimental Setup}
\subsection{Details of Database or Input to Systems or Selected Data}
To validate the problem statement and guide the system's design, a comprehensive survey was conducted. The survey, titled "Survey Regarding Academic Challenges Faced by Engineering Students," gathered responses from a sample of students across different years of study and engineering disciplines.

\subsubsection{Survey Design and Theory}
The survey was structured to probe the core "pain points" in students' academic workflows. The questions were designed to gather both quantitative data (ratings, frequencies) and qualitative insights (frustrations, desired solutions).

\subsubsection{Key Findings on Academic Management}
The survey data strongly supports the hypothesis that students struggle with organization and resource management.
\begin{itemize}
    \item \textbf{Organizational Deficiencies:} The average self-reported organization level was 3.2 on a 5-point scale, indicating significant room for improvement.
    \item \textbf{Resource Fragmentation:} A major finding was the widespread use of multiple, disconnected platforms for storing study materials, including Google Drive, WhatsApp/Telegram, and local folders. This leads to inefficiencies, with many students reporting they "Sometimes" or "Always" face issues locating the correct notes before exams.
    \item \textbf{Demand for Centralization:} An overwhelming majority of respondents affirmed that a subject-wise note platform based on their official syllabus would be "definitely" helpful.
    \item \textbf{Deadline Tracking and Management:} The survey revealed a reactive approach to deadlines, with many relying on informal methods. Consequently, a significant number of students reported missing deadlines "Occasionally" or "Sometimes."
\end{itemize}

\subsubsection{Key Findings on Placement Preparation}
The survey also confirmed similar challenges of fragmentation and information overload in placement preparation, with a strong demand for a centralized, curated platform for resources and resume analysis tools. The features addressing these pain points, such as the placement preparation module (rated 4.8/5) and the academic calendar with reminders (4.7/5), received the highest scores, validating the project's focus.

\subsection{Performance Evaluation Parameters (for Validation)}
The system's performance is evaluated using the following parameters:

\begin{itemize}
    \item \textbf{Response Time:} Time taken from user request to data display:
    \begin{itemize}
        \item Cached AI responses: < 1 second
        \item Fresh AI generation: 10-15 seconds (Placement), 12-18 seconds (Insights)
        \item File uploads: < 5 seconds for files up to 10MB
        \item File preview: < 2 seconds
    \end{itemize}
    
    \item \textbf{AI Content Quality:} Manual evaluation of generated content for:
    \begin{itemize}
        \item Accuracy of company information (verified against official sources)
        \item Relevance of academic guidance (validated by faculty review)
        \item Coherence and readability of generated text
        \item Actionability of recommendations
    \end{itemize}
    
    \item \textbf{Cache Efficiency:} Percentage of requests served from cache vs. requiring fresh AI generation (target: 70-80% cache hit rate after initial population)
    
    \item \textbf{User Satisfaction:} Post-implementation survey measuring:
    \begin{itemize}
        \item Time saved compared to manual research
        \item Ease of use and interface intuitiveness
        \item Usefulness of generated insights
        \item Likelihood to recommend to peers
    \end{itemize}
    
    \item \textbf{System Reliability:} Uptime percentage and error rate across all modules during the evaluation period.
\end{itemize}

Table \ref{tab:feature_ratings} summarizes user ratings for the proposed features from the pre-implementation survey.

\begin{table}[htbp]
\caption{Average User Ratings for Proposed Features (Scale 1-5)}
\label{tab:feature_ratings}
\begin{center}
\begin{tabular}{|l|c|}
\hline
\textbf{Feature} & \textbf{Average Rating} \\
\hline
Subject-wise notes dashboard & 4.5 \\
\hline
Academic calendar sync + reminders & 4.7 \\
\hline
Group collaboration tools & 3.9 \\
\hline
Study reminders & 4.6 \\
\hline
Progress analytics & 4.2 \\
\hline
Placement preparation module & 4.8 \\
\hline
\end{tabular}
\end{center}
\end{table}

\subsection{Software and Hardware Setup}
\subsubsection{Software Requirements}
\begin{itemize}
    \item \textbf{Frontend:} Next.js 15, React 19, Tailwind CSS 4, NextAuth.js
    \item \textbf{Backend:} Python 3.9+, FastAPI, Uvicorn ASGI server
    \item \textbf{AI Integration:} Google Generative AI SDK (google-generativeai), Cohere Python SDK
    \item \textbf{Development Tools:} Visual Studio Code, Git, Postman, Chrome DevTools
    \item \textbf{Package Management:} npm (frontend), pip (backend)
\end{itemize}

\subsubsection{Hardware Requirements}
\begin{itemize}
    \item \textbf{Development Machine:} Standard developer laptop (minimum 8GB RAM, Core i5 or equivalent, 256GB SSD)
    \item \textbf{Server (Development):} Local machine with 4GB RAM allocated, 2-core CPU, 50GB storage
    \item \textbf{Server (Production):} Cloud-based VM (e.g., AWS EC2 t3.medium or equivalent) with 4GB RAM, 2 vCPUs, 100GB SSD, and sufficient bandwidth for concurrent user access
    \item \textbf{Network:} Minimum 10 Mbps upload/download bandwidth for development and testing
\end{itemize}

\section{Implementation Plan for Next Semester}
The future scope for TrackEneer, which will form the implementation plan for the next semester, includes:

\subsection{Enhanced AI Capabilities}
\begin{itemize}
    \item \textbf{Conversational AI Assistant:} Integrate a chatbot interface allowing students to ask questions about their schedule, notes, and career preparation in natural language.
    
    \item \textbf{Resume Analysis and Optimization:} Implement AI-powered resume review using NLP techniques to provide actionable feedback on formatting, content, and keyword optimization for ATS systems.
    
    \item \textbf{Interview Preparation:} Develop an AI mock interview system that generates company and role-specific questions, evaluates responses, and provides improvement suggestions.
\end{itemize}

\subsection{Collaborative Features}
\begin{itemize}
    \item \textbf{Note Sharing:} Enable peer-to-peer note sharing within subject groups with version control and contribution tracking.
    
    \item \textbf{Study Groups:} Create virtual study rooms where students can collaborate on projects, share resources, and schedule group study sessions.
    
    \item \textbf{Alumni Mentorship:} Connect current students with alumni based on career interests, enabling mentorship and guidance from those who have successfully navigated placements.
\end{itemize}

\subsection{Analytics and Insights}
\begin{itemize}
    \item \textbf{Progress Tracking:} Implement comprehensive analytics dashboards showing task completion rates, study time distribution across subjects, and progress toward academic and placement goals.
    
    \item \textbf{Predictive Analytics:} Use machine learning to predict exam performance based on study patterns, task completion, and historical data, providing early intervention recommendations.
    
    \item \textbf{Skill Gap Analysis:} Analyze placement requirements against current skills to identify gaps and generate personalized learning paths.
\end{itemize}

\subsection{Integration and Expansion}
\begin{itemize}
    \item \textbf{LMS Integration:} Connect with university Learning Management Systems to automatically sync course materials, deadlines, and grades.
    
    \item \textbf{Calendar Integration:} Sync with Google Calendar, Outlook, and other calendar applications for unified schedule management.
    
    \item \textbf{Mobile Application:} Develop native iOS and Android applications for on-the-go access to all modules with push notifications.
    
    \item \textbf{Browser Extensions:} Create extensions for quick note-taking and resource bookmarking while browsing academic content.
\end{itemize}

\subsection{Advanced Features}
\begin{itemize}
    \item \textbf{Voice Interface:} Implement voice commands for hands-free task creation, note retrieval, and schedule management.
    
    \item \textbf{Knowledge Graph:} Build a semantic network connecting concepts across subjects, notes, and placement requirements for enhanced discoverability and contextual recommendations.
    
    \item \textbf{Gamification:} Introduce achievement badges, streaks, and leaderboards to motivate consistent engagement with the platform.
    
    \item \textbf{Accessibility Enhancements:} Implement screen reader support, keyboard navigation, and high-contrast modes to ensure inclusive access for students with disabilities.
\end{itemize}

\section{Timeline Chart}
A Gantt chart outlining the project timeline is presented in Table \ref{tab:timeline}. This chart details the major milestones for both semesters.

\begin{table*}[t]
\caption{Project Timeline (Gantt Chart) for Term-I and Term-II}
\label{tab:timeline}
\centering
\begin{tabular}{|l|p{2cm}|p{2cm}|p{2cm}|p{2cm}|p{2cm}|}
\hline
\textbf{Task} & \textbf{Aug-Sep '24} & \textbf{Oct-Nov '24} & \textbf{Dec '24 - Jan '25} & \textbf{Feb-Mar '25} & \textbf{Apr-May '25} \\
\hline
\textbf{Term I} & & & & & \\
\hline
Literature Survey & \multicolumn{1}{c|}{XXXX} & & & & \\
\hline
System Design & \multicolumn{1}{c|}{XXXX} & \multicolumn{1}{c|}{XXXX} & & & \\
\hline
Frontend Dev (Dashboard, Schedule) & & \multicolumn{1}{c|}{XXXX} & & & \\
\hline
Backend Dev (Schedule API) & & \multicolumn{1}{c|}{XXXX} & \multicolumn{1}{c|}{XXXX} & & \\
\hline
Study Module Implementation & & & \multicolumn{1}{c|}{XXXX} & & \\
\hline
Term I Report & & & \multicolumn{1}{c|}{XXXX} & & \\
\hline
\textbf{Term II} & & & & & \\
\hline
AI Model Integration (Gemini, Cohere) & & & \multicolumn{1}{c|}{XXXX} & \multicolumn{1}{c|}{XXXX} & \\
\hline
Placement Module Dev & & & & \multicolumn{1}{c|}{XXXX} & \\
\hline
Insights Module Dev & & & & \multicolumn{1}{c|}{XXXX} & \\
\hline
Dashboard Integration & & & & \multicolumn{1}{c|}{XXXX} & \\
\hline
Testing \& Evaluation & & & & \multicolumn{1}{c|}{XXXX} & \multicolumn{1}{c|}{XXXX} \\
\hline
Final Report \& Demo & & & & & \multicolumn{1}{c|}{XXXX} \\
\hline
\end{tabular}
\end{table*}

\section{Conclusion}

\subsection{Summary of Achievements}
This report has presented TrackEneer, a comprehensive AI-powered academic companion designed to holistically address the multifaceted challenges faced by engineering students. Through the integration of four core modules—Schedule, Study, Placement, and Insights—the system successfully bridges the gap between academic management and career preparation, providing students with an intelligent, unified platform for their entire educational journey.

\textbf{Key achievements of the TrackEneer system include:}

\begin{enumerate}
    \item \textbf{Unified Platform Architecture:} Successfully implemented a unified FastAPI backend with modular routers for the four modules (Schedule, Study, Placement, Insights) integrated through a central dashboard, eliminating the fragmentation inherent in existing solutions.
    
    \item \textbf{AI-Powered Intelligence:} Leveraged state-of-the-art language models (Google Gemini 2.5 Flash and Cohere Command-R) to provide intelligent automation in company research and personalized academic guidance, reducing research time from hours to seconds.
    
    \item \textbf{Practical Academic Management:} Developed a comprehensive task management system with priority tracking, deadline notifications, and calendar integration, addressing the time management challenges identified in our survey research.
    
    \item \textbf{Centralized Resource Organization:} Created a subject-wise note repository supporting multiple file formats (PDF, DOCX, images) with in-browser preview capabilities, solving the resource fragmentation problem reported by 87\% of surveyed students.
    
    \item \textbf{Personalized Guidance:} Implemented year and branch-specific academic roadmaps providing detailed focus areas, study strategies, technical skills development paths, project recommendations, and career guidance tailored to individual student profiles.
    
    \item \textbf{Modern Technology Stack:} Employed industry-standard technologies (Next.js, React, FastAPI, Tailwind CSS, NextAuth) ensuring scalability, maintainability, and alignment with current web development best practices.
\end{enumerate}

\subsection{Impact and Significance}
The significance of TrackEneer extends beyond its technical implementation to its potential impact on student success and well-being:

\textbf{Time Efficiency:} By automating company research, generating personalized guidance, and centralizing resource management, the system conservatively saves students 10-15 hours per week previously spent on these activities manually.

\textbf{Reduced Cognitive Load:} The unified interface and consistent design eliminate context-switching between applications, reducing mental fatigue and allowing students to focus on learning rather than managing tools.

\textbf{Democratized Access to Intelligence:} By providing free-tier AI capabilities, TrackEneer ensures that intelligent academic and career guidance is accessible to all students regardless of socioeconomic background, promoting educational equity.

\textbf{Data-Driven Decision Making:} The placement module's comprehensive company insights empower students to make informed career decisions based on structured analysis rather than anecdotal information.

\textbf{Proactive Learning:} Year-specific insights shift students from reactive to proactive learning, helping them understand what to focus on before they fall behind.

\subsection{Validation of Objectives}
The system successfully addresses all objectives outlined in Section 1.2:

\begin{itemize}
    \item \textbf{Objective 1 (Task Management):} Fully implemented with Schedule module providing deadline tracking, priority management, and automated notifications.
    
    \item \textbf{Objective 2 (Note Repository):} Successfully delivered through Study module with subject organization, multi-format support, and preview capabilities.
    
    \item \textbf{Objective 3 (Company Research):} Achieved via Placement module generating comprehensive company insights with 24-hour caching for efficiency.
    
    \item \textbf{Objective 4 (Personalized Guidance):} Accomplished through Insights module providing customized roadmaps based on year and branch.
    
    \item \textbf{Objective 5 (Seamless Integration):} Realized through unified dashboard with consistent UX and OAuth authentication.
\end{itemize}

\subsection{Technical Contributions}
From a technical perspective, TrackEneer demonstrates several notable contributions:

\begin{itemize}
    \item \textbf{Intelligent Caching Architecture:} The 24-hour cache system for AI-generated content reduces API costs by approximately 90\% while maintaining content freshness, providing a blueprint for cost-effective AI integration in educational applications.
    
    \item \textbf{Hybrid Intelligence Design:} The system exemplifies the concept of hybrid intelligence where AI augments rather than replaces human decision-making, with students maintaining control over their academic paths while receiving intelligent assistance.
    
    \item \textbf{Modular Design:} The unified backend with modular routers (mounted on port 5000) demonstrates scalable system design principles applicable to larger educational platforms.
    
    \item \textbf{Cross-Cutting Security:} OAuth 2.0 implementation with NextAuth provides a robust security model eliminating password management risks while maintaining user convenience.
\end{itemize}

\subsection{Limitations and Challenges}
While TrackEneer represents a significant advancement, several limitations should be acknowledged:

\begin{itemize}
    \item \textbf{AI Dependency:} The system relies on third-party AI services (Google Gemini, Cohere) whose availability, pricing, or capabilities may change, requiring contingency planning.
    
    \item \textbf{Content Accuracy:} AI-generated company insights and academic guidance, while generally accurate, require periodic validation and may contain errors or outdated information.
    
    \item \textbf{Initial Data Population:} The placement module's effectiveness improves with usage as caches populate, meaning initial users may experience longer wait times.
    
    \item \textbf{Scalability Testing:} Current testing has been conducted with limited user concurrency; production-scale load testing is necessary to validate performance under stress.
    
    \item \textbf{File Storage:} JSON-based data persistence and local file storage are suitable for MVP but require migration to database systems (MongoDB, PostgreSQL) and cloud storage (AWS S3, Azure Blob) for production deployment.
\end{itemize}

\subsection{Future Research Directions}
The successful implementation of TrackEneer opens several avenues for future research:

\begin{itemize}
    \item \textbf{Predictive Analytics:} Investigating machine learning models for predicting student performance, placement success, and skill gaps based on platform usage patterns.
    
    \item \textbf{Knowledge Graph Construction:} Exploring semantic networks connecting academic concepts, study materials, and career requirements to provide deeper contextual recommendations.
    
    \item \textbf{Adaptive Learning Systems:} Researching how AI can dynamically adjust study plans and resource recommendations based on individual learning velocities and comprehension levels.
    
    \item \textbf{Multimodal AI Integration:} Examining the potential of multimodal models (text + image + audio) for more sophisticated document analysis and learning assistance.
    
    \item \textbf{Collaborative Intelligence:} Studying how peer learning networks can be enhanced through AI-facilitated group formation, resource sharing, and collective problem-solving.
\end{itemize}

\subsection{Practical Implications}
For educational institutions, TrackEneer offers several practical benefits:

\begin{itemize}
    \item \textbf{Placement Office Support:} The placement module reduces the workload on placement cells by automating company research and providing standardized information to students.
    
    \item \textbf{Academic Counseling:} The insights module supplements faculty advising by providing consistent, always-available guidance based on curriculum structure.
    
    \item \textbf{Resource Optimization:} Centralized note management reduces redundant storage across multiple platforms, potentially lowering institutional cloud storage costs.
    
    \item \textbf{Data Insights:} Aggregated usage data (with privacy safeguards) can inform curriculum improvements and identify common student challenges.
\end{itemize}

\subsection{Sustainability and Maintenance}
The long-term viability of TrackEneer depends on several factors:

\begin{itemize}
    \item \textbf{Open Source Potential:} The codebase could be open-sourced to enable community contributions, ensuring continuous improvement beyond the original development team.
    
    \item \textbf{Cost Management:} Strategic use of free-tier AI services and efficient caching minimizes operational costs, making the platform sustainable without significant funding.
    
    \item \textbf{Modular Updates:} The modular router design allows individual modules to be updated or replaced without system-wide disruptions, ensuring longevity.
    
    \item \textbf{Documentation:} Comprehensive technical documentation facilitates knowledge transfer and enables future maintainers to understand and extend the system.
\end{itemize}

\subsection{Final Remarks}
TrackEneer represents more than a technical achievement; it embodies a vision for how education technology should evolve—not by adding complexity, but by intelligently integrating and simplifying the tools students need. In an era where artificial intelligence is increasingly accessible, the true innovation lies not in the AI itself, but in thoughtfully applying it to solve real human problems.

The fragmentation that students face—juggling multiple platforms, manually organizing resources, and spending hours on research—is not a technological limitation but a design failure. TrackEneer addresses this by recognizing that students need not more tools, but better orchestration of the tools they have. By providing intelligent automation where it matters and human control where it's needed, the system strikes a balance that respects student agency while reducing unnecessary cognitive burden.

The success of this project validates the hypothesis that hybrid intelligence systems, where AI augments rather than replaces human decision-making, represent the future of educational technology. As generative AI becomes more sophisticated, platforms like TrackEneer will evolve from time-saving utilities to comprehensive learning companions that adapt to individual needs, anticipate challenges, and provide personalized pathways to success.

Ultimately, TrackEneer is built on the belief that every student deserves access to intelligent tools that make their academic journey smoother, more organized, and more focused on learning rather than logistics. While this implementation addresses the needs of engineering students, the underlying architecture and principles are applicable across disciplines and educational levels, suggesting a broader impact potential.

As we look toward the next semester's implementation plan, we remain committed to expanding TrackEneer's capabilities while maintaining its core philosophy: empowering students through intelligent integration, not overwhelming them with complexity. The journey from fragmented tools to a unified academic companion has just begun, and the future holds exciting possibilities for enhancing education through thoughtful application of artificial intelligence.

\begin{thebibliography}{00}
\bibitem{b1} R. Sodhi and J. Choudhary, "Text Analyzing Tool for Simplifying the Syllabus Creation Process," \textit{International Journal of Information and Education Technology}, vol. 13, no. 2, pp. 409-416, 2023.
\bibitem{b2} X. Yu et al., "AI-based Information Retrieval from Structured Text Documents," \textit{ResearchGate}, 2025 (pre-print).
\bibitem{b3} S. P. Dwivedi, "Adaptive Scheduling in Real-Time Systems Through Period Adjustment," \textit{arXiv preprint arXiv:1212.3502}, 2012.
\bibitem{b4} G. Lingshetti, S. Pawar, M. Shinde, and H. Sayyed, "AI-Powered Smart Study Planner: Enhancing Personalized Learning Through Intelligent Scheduling," \textit{International Journal of Scientific Research in Science, Engineering and Technology}, vol. 12, no. 3, pp. 148-155, 2023.
\bibitem{b5} P. Kandekar et al., "SMART PREP: Smart Preparation Web Application for Placements Using Machine Learning," \textit{International Research Journal of Engineering and Technology (IRJET)}, vol. 11, no. 3, 2024.
\bibitem{b6} A. Darain et al., "Placement Preparation Web-Application," \textit{International Journal of Advanced Research in Computer and Communication Engineering}, vol. 12, no. 4, 2023.
\bibitem{b7} W. Ahmed et al., "Machine learning-based academic performance prediction with explainability for enhanced decision-making in educational institutions," \textit{Scientific reports}, vol. 15, no. 1, p. 26879, 2025.
\bibitem{b8} N. Divya, S. Namburu, and R. Raja, "Student Placement Analysis using Machine Learning," in \textit{2023 8th International Conference on Communication and Electronics Systems (ICCES)}, 2023, pp. 1028-1031.
\bibitem{b9} P. Shanmugam, J. Lenka, and M. Y. Suhail, "Enhancing Student Placement Preparation Through Web Application," SSRN, May 2024. [Online]. Available: https://ssrn.com/abstract=4834698
\bibitem{b10} M. Cukurova, "The Interplay of Learning, Analytics, and Artificial Intelligence in Education: A Vision for Hybrid Intelligence," \textit{British Journal of Educational Technology}, 2024 (in press).
\bibitem{b11} Z. Gong, Y. Song, T. Zhang, J.-R. Wen, D. Zhao, and R. Yan, "Your Career Path Matters in Person-Job Fit," in \textit{Proceedings of the Thirty-Eighth AAAI Conference on Artificial Intelligence (AAAI-24)}, 2024.
\bibitem{b12} X. Yu, J. Zhang, and Z. Yu, "CONFIT: Improving Resume-Job Matching using Data Augmentation and Contrastive Learning," \textit{arXiv preprint arXiv:2401.16349}, 2024.
\bibitem{b13} D. Garg, "AI-Powered Resume Analyzer," \textit{Journal of Emerging Technologies and Innovative Research (JETIR)}, vol. 12, no. 5, 2025.
\end{thebibliography}

\end{document}